%
% menu_lifecycle_d1_buggy.tex
%
% FIDELITY CONTROL for docs/menu_lifecycle.tex, defect D1.
%
% This spec differs from the fixed menu_lifecycle.tex in exactly one place:
% MenuDeliver (the atomic get-then-put, the D1 fix) is replaced by the
% shipped code's two-phase offer -- OfferBegin captures the queue reference
% under _inboxes_lock, the lock is RELEASED, then OfferEnd puts and returns
% "delivered" unconditionally. A DepartGraceful interleaving between the two
% pops the queue, so OfferEnd's put lands in an orphan while still reporting
% "delivered" -- the false-positive delivery M1 forbids.
%
% Expected: M1's negation is reachable.
%   probcli docs/menu_lifecycle_d1_buggy.tex -model_check -nodead \
%     -goal "reported = ddelivered & lost = fset" -p DEFAULT_SETSIZE 3
%   => goal FOUND (a witness through OfferBegin; DepartGraceful; OfferEnd).
%
% Against the fixed spec the same goal is NOT found. This control isolates
% D1: it leaves SetMenu atomic (D3 fixed) and DepartLapsed a full sweep
% (D2 fixed), so the only reachable violation is the offer window.
%

\documentclass[a4paper,10pt,fleqn]{article}
\usepackage[margin=1in]{geometry}
\usepackage{fuzz}
\usepackage[colorlinks=true,linkcolor=blue,citecolor=blue,urlcolor=blue]{hyperref}

\begin{document}

\title{lux Agent-Menu Lifecycle --- D1 Fidelity Control}
\author{The non-atomic offer: get-then-put across a lock release}
\date{September 2026}
\maketitle

This is a fidelity control for \texttt{menu\_lifecycle.tex}. Only the
delivery operation differs; every other schema is copied verbatim.

\begin{zed}
  [CONN]
\end{zed}

\begin{zed}
  DSTATE ::= dnone | ddelivered | dgone
\end{zed}

\begin{zed}
  FSTATE ::= fclear | fset
\end{zed}

\begin{zed}
  APHASE ::= aidle | astore
\end{zed}

\begin{zed}
  OPHASE ::= oidle | oput
\end{zed}

\begin{schema}{MenuReg}
  registered : \power CONN \\
  identified : \power CONN \\
  lapsed : \power CONN \\
  hasInbox : \power CONN \\
  menuOwner : \power CONN \\
  reported : DSTATE \\
  lost : FSTATE \\
  admitPhase : APHASE \\
  admitFor : \power CONN \\
  offerPhase : OPHASE \\
  offerFor : \power CONN
\where
  identified \subseteq registered \\
  lapsed \subseteq registered \\
  \# admitFor \leq 1 \\
  \# offerFor \leq 1
\end{schema}

\begin{schema}{Init}
  MenuReg'
\where
  registered' = \emptyset \\
  identified' = \emptyset \\
  lapsed' = \emptyset \\
  hasInbox' = \emptyset \\
  menuOwner' = \emptyset \\
  reported' = dnone \\
  lost' = fclear \\
  admitPhase' = aidle \\
  admitFor' = \emptyset \\
  offerPhase' = oidle \\
  offerFor' = \emptyset
\end{schema}

\begin{schema}{Connect}
  \Delta MenuReg \\
  c? : CONN
\where
  registered' = registered \cup \{ c? \} \\
  lapsed' = lapsed \setminus \{ c? \} \\
  identified' = identified \\
  hasInbox' = hasInbox \\
  menuOwner' = menuOwner \\
  reported' = reported \\
  lost' = lost \\
  admitPhase' = admitPhase \\
  admitFor' = admitFor \\
  offerPhase' = offerPhase \\
  offerFor' = offerFor
\end{schema}

\begin{schema}{Identify}
  \Delta MenuReg \\
  c? : CONN
\where
  c? \in registered \\
  identified' = identified \cup \{ c? \} \\
  registered' = registered \\
  lapsed' = lapsed \\
  hasInbox' = hasInbox \\
  menuOwner' = menuOwner \\
  reported' = reported \\
  lost' = lost \\
  admitPhase' = admitPhase \\
  admitFor' = admitFor \\
  offerPhase' = offerPhase \\
  offerFor' = offerFor
\end{schema}

\begin{schema}{Arm}
  \Delta MenuReg \\
  c? : CONN
\where
  c? \in registered \\
  hasInbox' = hasInbox \cup \{ c? \} \\
  registered' = registered \\
  identified' = identified \\
  lapsed' = lapsed \\
  menuOwner' = menuOwner \\
  reported' = reported \\
  lost' = lost \\
  admitPhase' = admitPhase \\
  admitFor' = admitFor \\
  offerPhase' = offerPhase \\
  offerFor' = offerFor
\end{schema}

\begin{schema}{LeaseLapse}
  \Delta MenuReg \\
  c? : CONN
\where
  c? \in registered \\
  c? \notin lapsed \\
  lapsed' = lapsed \cup \{ c? \} \\
  registered' = registered \\
  identified' = identified \\
  hasInbox' = hasInbox \\
  menuOwner' = menuOwner \\
  reported' = reported \\
  lost' = lost \\
  admitPhase' = admitPhase \\
  admitFor' = admitFor \\
  offerPhase' = offerPhase \\
  offerFor' = offerFor
\end{schema}

\begin{schema}{SetMenu}
  \Delta MenuReg \\
  c? : CONN
\where
  c? \in registered \\
  c? \in identified \\
  lapsed' = lapsed \setminus \{ c? \} \\
  hasInbox' = hasInbox \cup \{ c? \} \\
  menuOwner' = menuOwner \cup \{ c? \} \\
  registered' = registered \\
  identified' = identified \\
  reported' = reported \\
  lost' = lost \\
  admitPhase' = admitPhase \\
  admitFor' = admitFor \\
  offerPhase' = offerPhase \\
  offerFor' = offerFor
\end{schema}

\begin{schema}{SetMenuRefused}
  \Delta MenuReg \\
  c? : CONN
\where
  c? \notin identified \\
  lapsed' = lapsed \setminus \{ c? \} \\
  registered' = registered \\
  identified' = identified \\
  hasInbox' = hasInbox \\
  menuOwner' = menuOwner \\
  reported' = reported \\
  lost' = lost \\
  admitPhase' = admitPhase \\
  admitFor' = admitFor \\
  offerPhase' = offerPhase \\
  offerFor' = offerFor
\end{schema}

\begin{schema}{DepartGraceful}
  \Delta MenuReg \\
  c? : CONN
\where
  c? \in registered \\
  registered' = registered \setminus \{ c? \} \\
  identified' = identified \setminus \{ c? \} \\
  lapsed' = lapsed \setminus \{ c? \} \\
  hasInbox' = hasInbox \setminus \{ c? \} \\
  menuOwner' = menuOwner \setminus \{ c? \} \\
  reported' = reported \\
  lost' = lost \\
  admitPhase' = admitPhase \\
  admitFor' = admitFor \\
  offerPhase' = offerPhase \\
  offerFor' = offerFor
\end{schema}

\begin{schema}{DepartLapsed}
  \Delta MenuReg \\
  c? : CONN
\where
  c? \in registered \\
  c? \in lapsed \\
  registered' = registered \setminus \{ c? \} \\
  identified' = identified \setminus \{ c? \} \\
  lapsed' = lapsed \setminus \{ c? \} \\
  hasInbox' = hasInbox \setminus \{ c? \} \\
  menuOwner' = menuOwner \setminus \{ c? \} \\
  reported' = reported \\
  lost' = lost \\
  admitPhase' = admitPhase \\
  admitFor' = admitFor \\
  offerPhase' = offerPhase \\
  offerFor' = offerFor
\end{schema}

% -- THE DEFECT: offer's get and put split across a lock release ---------

% OfferBeginPresent models offer reading _inboxes.get(c) under the lock,
% finding a queue, and RELEASING the lock while holding the reference --
% the connection is captured into offerFor and the put is now committed.
\begin{schema}{OfferBeginPresent}
  \Delta MenuReg \\
  c? : CONN
\where
  offerPhase = oidle \\
  c? \in registered \\
  c? \in hasInbox \\
  offerPhase' = oput \\
  offerFor' = \{ c? \} \\
  registered' = registered \\
  identified' = identified \\
  lapsed' = lapsed \\
  hasInbox' = hasInbox \\
  menuOwner' = menuOwner \\
  reported' = reported \\
  lost' = lost \\
  admitPhase' = admitPhase \\
  admitFor' = admitFor
\end{schema}

% OfferBeginAbsent models offer's get returning None (no inbox): the honest
% provider_gone, no window opened.
\begin{schema}{OfferBeginAbsent}
  \Delta MenuReg \\
  c? : CONN
\where
  offerPhase = oidle \\
  c? \in registered \\
  c? \notin hasInbox \\
  reported' = dgone \\
  registered' = registered \\
  identified' = identified \\
  lapsed' = lapsed \\
  hasInbox' = hasInbox \\
  menuOwner' = menuOwner \\
  lost' = lost \\
  admitPhase' = admitPhase \\
  admitFor' = admitFor \\
  offerPhase' = offerPhase \\
  offerFor' = offerFor
\end{schema}

% OfferEnd models put-then-return-True. It reports ddelivered
% UNCONDITIONALLY -- the bug. If the captured connection has since lost its
% inbox (a drop_session ran in the window), the put lands in an orphan and
% lost becomes fset, yet the caller is still told "delivered".
\begin{schema}{OfferEnd}
  \Delta MenuReg
\where
  offerPhase = oput \\
  reported' = ddelivered \\
  (offerFor \subseteq hasInbox \implies lost' = lost) \\
  (\lnot (offerFor \subseteq hasInbox) \implies lost' = fset) \\
  offerPhase' = oidle \\
  offerFor' = \emptyset \\
  registered' = registered \\
  identified' = identified \\
  lapsed' = lapsed \\
  hasInbox' = hasInbox \\
  menuOwner' = menuOwner \\
  admitPhase' = admitPhase \\
  admitFor' = admitFor
\end{schema}

% MenuProviderGone still models the router's live-set read failing.
\begin{schema}{MenuProviderGone}
  \Delta MenuReg \\
  c? : CONN
\where
  c? \notin registered \\
  reported' = dgone \\
  registered' = registered \\
  identified' = identified \\
  lapsed' = lapsed \\
  hasInbox' = hasInbox \\
  menuOwner' = menuOwner \\
  lost' = lost \\
  admitPhase' = admitPhase \\
  admitFor' = admitFor \\
  offerPhase' = offerPhase \\
  offerFor' = offerFor
\end{schema}

\end{document}
